\section*{Notation}\label{sec:notation}
\phantomsection\addcontentsline{toc}{section}{Notation}
This registry records the symbols whose meaning is not uniform across
the compendium. A symbol listed here with a section scope is local to
that section; the canonical (unqualified) form is the fixed-map
propagator defined in the front matter. Renames applied in Pass 1 are
mechanical token replacements, listed at the end of this registry; no
statement was reworded by a rename. Symbols not listed here are used
with their standard meaning or are defined where they first appear.

\begin{center}
\small
\begin{tabular}{>{\raggedright\arraybackslash}p{0.18\textwidth}>{\raggedright\arraybackslash}p{0.46\textwidth}>{\raggedright\arraybackslash}p{0.24\textwidth}}
Symbol & Meaning & Scope / location \\
\hline
\(U_N\) & Fixed kicked propagator \(U_N=\exp(i\tau N A_N)M_N\),
  \(\tau=10^{-80}\) fixed & Canonical; front matter, fixed-map sections \\
\(U_{N,\mathrm{loc}}\) & A cat quantization with the exact covariance
  \(UQ_N(v)U^*=Q_N(Tv)\) & Locality section (Ch.~9), renamed in Pass 1 \\
\(U_{N,\mathrm{rig}}\) & A cat quantization taken with the exact
  covariance, phase immaterial & Resonance-rigidity section (Ch.~10),
  renamed in Pass 1 \\
\(U_{N,\mathrm{Fl}}\) & The Floquet operator \(U_N=VU_0V^*\) of the
  finite Floquet construction & Floquet section (Ch.~7), renamed in Pass 1 \\
\(U_{N,\mathrm{test}}\) & A generic test unitary (any unitary
  quantization satisfying the stated Egorov estimate) &
  Spectral-resolution section (Ch.~11), renamed in Pass 1 \\
\(U_N(T)\) & The \(T\)-parametrized family \(U_N(T)=VS_h^{-1}V^*\) &
  Pell-blocks section; the argument distinguishes it from \(U_N\) \\
\(\mathcal{M}_r\) & \(\sum_v|\widehat a(v)|(1+|v|_1)^r\) &
  Locality section \\
\(M_r\) & \(\sum_v|\widehat a(v)||v|_2^r\) (Euclidean frequency weight) &
  Return-words section \\
\(M_j\) & The same object as \(M_r\), with index letter \(j\) &
  Analytic section (regular-gradient rigidity) \\
\(M_r^{\mathrm{mom}}\), \(M_2^{\mathrm{mom}}\), \(M_4^{\mathrm{mom}}\) &
  Moments \(\mathbb{E}|Z|^r\) of \(Z=s_N(w)\) & Fixed-map frontier,
  early moment block, renamed in Pass 1 \\
\(M_s\) & Occupancy moments \(N^{-1}\sum_h n_h^s\) &
  Fixed-map frontier, resonant-cubic section (defined in place) \\
\(M_3, M_4\) & Occupancy moments \(M_s\) at \(s=3,4\) &
  Fixed-map frontier, resonant-cubic and folded-square sections
  (verified Pass 2: same \(n_h\) occupancies and \(E_N\) as the
  \(M_s\) definition) \\
\(\Phi^t\) & Flow of \(X_g\) (locality convention) &
  Locality section \\
\(\Phi_g^t\) & Flow of \(X_g\) (same convention, explicit generator) &
  Locality section \\
\(\Psi_h^t\) & Flow of \(X_h=(h_p,-h_q)\); inverse of the locality flow &
  Return-words section, renamed from \(\Phi_h^t\) in Pass 1 \\
\(\kappa_0\) & Rescaled time; locality writes
  \(e^{i\kappa_0N\operatorname{Op}_N(f)}\), return words write
  \(\kappa=\kappa_0/(2\pi)\) & Locality section \\
\(\kappa\) & Classical time itself, through \(e=\kappa/\hbar_N\) &
  Return-words section \\
\(\kappa_0^{\mathrm{thr}}\) & The threshold constant
  \(\frac43(\sqrt{101}\sin(\pi/7)-1)/100\) & Fixed-map frontier,
  renamed in Pass 1 \\
\(\kappa\) & A local exponent or contraction parameter (e.g.\
  \(m=\lceil\ell^{1+\kappa}\rceil\); \(\kappa=1-\rho\)) &
  Section-local; defined where used \\
\(\kappa_4\) & Fourth cumulant & Standard notation, wherever used \\
\end{tabular}
\end{center}

\noindent\textbf{Pass 1 renames (mechanical, token-for-token).}
Fixed-map frontier, moment block (lines 378--414 only):
\(M_r, M_2, M_4 \to M_r^{\mathrm{mom}}, M_2^{\mathrm{mom}},
M_4^{\mathrm{mom}}\) (11 replacements).
Fixed-map frontier, threshold block (lines 1679--1681):
\(\kappa_0 \to \kappa_0^{\mathrm{thr}}\) (3 replacements).
Return words: \(\Phi_h^t \to \Psi_h^t\) (5 replacements).
Locality: \(U_N \to U_{N,\mathrm{loc}}\) (15 replacements).
Resonance rigidity: \(U_N \to U_{N,\mathrm{rig}}\) (9 replacements).
Floquet: \(U_N \to U_{N,\mathrm{Fl}}\) (3 replacements), plus the
coboundary cross-reference to the Floquet operator (1 replacement).
Spectral resolution: \(U_N \to U_{N,\mathrm{test}}\) (10 replacements)
and \(\widetilde U_N \to \widetilde{U_{N,\mathrm{test}}}\) (3 replacements).
The tags were placed in the subscript (\(U_{N,\mathrm{loc}}\) etc.) so
that existing superscripts (\(U_N^*, U_N^j, U_N^{-j}, U_N^K\)) keep
working; a superscript tag was tried first and broke the build with
double superscripts.

\noindent\textbf{Open items and Pass 2 resolutions.}
(1) RESOLVED in Pass 2: Fixed-map frontier, late \(M_3, M_4\)
(lines 1917, 1944, 1964 in the resonant-cubic section itself, and
line 2007 in the folded-square section, which reuses TRACE-011's
\(E_N\) and counterterm budget and contrasts the \(M_4\) bound with
the new \(G_d\) bound) are the occupancy moments \(M_s\) defined at
line 1910; see the \(M_3, M_4\) row above. No token rename was
applied. (2) \(\kappa\) as a section-local parameter is
documented above but not unified; unification would require
per-section renaming with proof reading, which is Pass 2 work.
(2) DECIDED in Pass 2 (no rename): the microscopic-paths \(\kappa>0\)
(free exponent parameter, e.g.\ \(m=\lceil\ell^{1+\kappa}\rceil\)) and
the unmasked \(\kappa=1-\rho>0\) are section-local parameters with
distinct definitions; renaming either would churn many formulas for no
mathematical benefit, and the registry entries above already
disambiguate them. Both stay as written.
(3) The analytic \(M_j\) and return-words \(M_r\) are the same object
under different index letters; no rename was applied, since the
registry entry already disambiguates them.
